\chapter{Verifikation, Validierung und Traceability}

\section{Verifikationsprinzip}

Ein Systemmodell ist nur dann nuetzlich, wenn es bis zum Nachweis fuehrt. Verifikation prueft, ob das System die spezifizierten Anforderungen erfuellt. Validierung prueft, ob das entwickelte System das tatsaechliche Problem der Stakeholder loest. Diese Unterscheidung entspricht dem im Lehrbuch beschriebenen Verifikations- und Validierungsverstaendnis \cite[Kap.~19]{becker2026sysml}.

\section{Initiale Testfaelle}

\begin{longtable}{L{0.14\textwidth} L{0.24\textwidth} L{0.22\textwidth} L{0.24\textwidth}}
\toprule
ID & Testfall & Verifiziert & Erwartetes Ergebnis \\
\midrule
\modeltable{tests_a1}
\bottomrule
\end{longtable}

Alle Testfaelle sind als Verifikationsdefinitionen im \sysml-Modell (\texttt{06\_tests.sysml}) hinterlegt. Keiner ist bisher an einem realen Roboter oder in einem dokumentierten Simulationslauf ausgefuehrt worden.

\section{Simulationspruefstand und Modellpruefung}

Gemaess DEC-014 wird zunaechst ausschliesslich in der Simulation verifiziert. Das Modell enthaelt dafuer einen logischen Pruefstand (\texttt{SimulationTestBench}) aus Roboter, Simulationswelt, Steuerungsadapter und Beobachter. Simulator, Weltgeometrie, Sensor- und Energiemodell sind noch nicht ausgewaehlt.

TC-NAV-002 und TC-UI-003 besitzen auswertbare Kriterien. Ohne vollstaendige Nachweise lautet das Ergebnis \enquote{inconclusive}, nicht \enquote{pass}. Die Modellpruefung (\texttt{tools/validate\_sysml.py}) laedt das Gesamtmodell mit der offiziellen SysML-v2-Referenzimplementierung, wertet synthetische Testdaten aus und prueft absichtlich beschaedigte Modellkopien. Synthetische Testdaten sind keine Beobachtungen.

\section{Testfallstruktur}

Jeder Testfall wird spaeter in folgender Struktur beschrieben:
\begin{itemize}
  \item Test-ID,
  \item verifizierte Requirement-ID,
  \item Ziel,
  \item Vorbedingungen,
  \item Testaufbau oder Simulationsumgebung,
  \item Stimulus oder Ablauf,
  \item erwartetes Ergebnis,
  \item Bewertungskriterium,
  \item benoetigte Messdaten,
  \item Status und letzter Durchlauf.
\end{itemize}

\section{Traceability-Kette}

Die minimale Traceability-Kette fuer dieses Projekt lautet:
\begin{quote}
Stakeholder Need $\rightarrow$ Requirement $\rightarrow$ Use Case oder Funktion $\rightarrow$ logisches Systemelement $\rightarrow$ Schnittstelle oder \ros-Element $\rightarrow$ Testfall.
\end{quote}

\begin{figure}[htbp]
  \centering
  \resizebox{\textwidth}{!}{\begin{tikzpicture}[
  item/.style={draw, rounded corners=2pt, align=center, minimum width=26mm, minimum height=8mm, fill=gray!8},
  arrow/.style={-{Latex[length=2mm]}, thick}
]
\node[item] (need) {Stakeholder\\Need};
\node[item, right=7mm of need] (req) {Requirement};
\node[item, right=7mm of req] (uc) {Use Case\\Betriebsfall};
\node[item, right=7mm of uc] (fun) {Funktion\\Verhalten};
\node[item, right=7mm of fun] (arch) {Systemelement\\Schnittstelle};
\node[item, right=7mm of arch] (test) {Test\\Nachweis};

\draw[arrow] (need) -- (req);
\draw[arrow] (req) -- (uc);
\draw[arrow] (uc) -- (fun);
\draw[arrow] (fun) -- (arch);
\draw[arrow] (arch) -- (test);
\end{tikzpicture}
}
  \caption{Minimale Traceability-Kette der Dokumentation}
  \label{fig:traceability-chain}
\end{figure}

\section{Traceability-Regeln}

\begin{itemize}
  \item Keine MVP-Anforderung ohne Testfall.
  \item Kein \ros-Element ohne modellierte Verantwortung.
  \item Keine Schnittstelle ohne Producer, Consumer und Zweck.
  \item Keine Architekturentscheidung ohne betrachtete Alternative.
  \item Jede geaenderte Anforderung loest eine Impact-Analyse auf Funktionen, Schnittstellen, Tests und Dokumentation aus.
\end{itemize}
